\section{Method}
\label{sec:method}

We present \textbf{AGILE}, an adaptive generalization framework that expands the
training distribution of a Vision-Language-Action (VLA) policy according to the
policy's measured competence. Instead of training on a fixed set of task
configurations, AGILE organizes the configuration space into progressively
larger regions and automatically switches to the next region when the current
policy is sufficiently reliable. The key idea is to use a world-model-induced
representation to estimate relations among task configurations, while using
online policy success as feedback for when to expand the training distribution.

\subsection{Problem Formulation}
\label{subsec:problem_formulation}

We consider a robotic manipulation task family parameterized by a task
configuration $s \in \mathcal{S}$. In our setting, a configuration specifies
factors such as object initial poses, camera pose, and visual appearance of the
workspace. Given an observation $o_t$, language instruction $l$, and robot state,
a VLA policy $\pi_\theta$ predicts an action chunk
$a_{t:t+H-1} \sim \pi_\theta(\cdot \mid o_t, l)$. The goal is to learn a policy
that succeeds not only on a narrow training distribution, but across a broader
configuration set $\mathcal{S}$.

Directly training on the full configuration set can be inefficient because many
configurations are initially too difficult for the policy, producing sparse or
low-quality learning signals. Conversely, training on a small fixed subset
limits generalization. AGILE addresses this by maintaining a stage-dependent
training set $\mathcal{A}_k \subseteq \mathcal{S}$, where $k$ is the current
expansion stage and
\begin{equation}
    \mathcal{A}_0 \subset \mathcal{A}_1 \subset \cdots \subset \mathcal{A}_K.
\end{equation}
The training distribution at stage $k$ samples configurations from
$\mathcal{A}_k$. Stage transitions are triggered automatically by policy
performance rather than by a fixed hand-designed training schedule.

\subsection{World-Model Representation of Task Configurations}
\label{subsec:world_model_representation}

AGILE assumes access to a world model or representation model that maps each
task configuration $s_i$ to an embedding vector $z_i \in \mathbb{R}^d$:
\begin{equation}
    z_i = f_{\mathrm{WM}}(s_i).
\end{equation}
The embedding is used to estimate how close a candidate configuration is to a
reference or already mastered region. We use cosine distance as the default
distance measure:
\begin{equation}
    d(i,j) = 1 - \frac{z_i^\top z_j}{\|z_i\|_2 \|z_j\|_2}.
\end{equation}
This distance is not treated as an absolute task difficulty label. Instead, it
provides a structured ordering over task configurations, allowing AGILE to
expand training from familiar regions to less familiar ones.

\subsection{Stage Construction}
\label{subsec:stage_construction}

AGILE supports two ways to construct expansion stages.

\paragraph{Reference-distance expansion.}
Given a reference configuration $s_{\mathrm{ref}}$, we sort all candidate
configurations by their distance to $s_{\mathrm{ref}}$ and choose cumulative
stage sets according to predefined stage sizes:
\begin{equation}
    \mathcal{A}_k =
    \{s_i \mid \operatorname{rank}_{d(i,\mathrm{ref})}(s_i) \leq n_k\},
\end{equation}
where $n_k$ is the cumulative number of configurations allowed at stage $k$.
This produces a simple and reproducible expansion from the reference region.

\paragraph{Dynamic region growing.}
To reduce reliance on a single reference point, we also consider a dynamic
top-$m$ expansion strategy. Let $\mathcal{M}_{k-1}$ denote the mastered region
before constructing stage $k$, initialized by either a reference configuration
or a base set. For each unselected candidate $s_i$, AGILE computes its average
distance to the closest $m$ configurations in $\mathcal{M}_{k-1}$:
\begin{equation}
    D(i, \mathcal{M}_{k-1}) =
    \frac{1}{m}
    \sum_{s_j \in \mathrm{NN}_m(s_i, \mathcal{M}_{k-1})} d(i,j).
\end{equation}
The next stage adds the candidates with the smallest values of
$D(i, \mathcal{M}_{k-1})$:
\begin{equation}
    \Delta_k =
    \operatorname{TopSmallest}_{n_k - n_{k-1}}
    \{D(i, \mathcal{M}_{k-1}) \mid s_i \notin \mathcal{A}_{k-1}\},
\end{equation}
\begin{equation}
    \mathcal{A}_k = \mathcal{A}_{k-1} \cup \Delta_k.
\end{equation}
This creates a data-driven expansion path over the task-configuration space.
Each stage is stored as a manifest containing both the newly added
configurations $\Delta_k$ and the cumulative allowed set $\mathcal{A}_k$.

\subsection{Adaptive Distribution Expansion During Training}
\label{subsec:adaptive_expansion}

During policy learning, AGILE samples environment resets only from the current
allowed set $\mathcal{A}_k$. Let $r_i \in \{0,1\}$ denote whether an episode on
configuration $s_i$ succeeds. At training step $t$, we estimate the success rate
over rollouts collected from the current stage:
\begin{equation}
    \hat{p}_t =
    \frac{1}{|\mathcal{B}_t|}
    \sum_{(s_i,r_i)\in\mathcal{B}_t} r_i,
\end{equation}
where $\mathcal{B}_t$ is the set of evaluated rollouts in the current update.
The stage is advanced when the measured competence exceeds a threshold
$\tau_k$ for a required number of hits:
\begin{equation}
    \sum_{t' \in \mathcal{H}_k}
    \mathbf{1}[\hat{p}_{t'} \geq \tau_k]
    \geq h_k,
\end{equation}
where $\mathcal{H}_k$ stores observations within the current stage and $h_k$ is
the required number of successful hits. In our implementation, hits do not need
to be consecutive unless explicitly configured. This avoids overreacting to a
single noisy evaluation step while still allowing the training distribution to
expand automatically once the policy becomes reliable.

We further support a grouped moving-average criterion that separates the
performance on previously available configurations and newly added
configurations. For stage $k>0$, let $\mathcal{A}_{k-1}$ be the old set and
$\Delta_k$ be the new increment. We compute:
\begin{equation}
    \hat{p}^{\mathrm{old}}_t =
    \mathbb{E}[r_i \mid s_i \in \mathcal{A}_{k-1}],
    \quad
    \hat{p}^{\mathrm{new}}_t =
    \mathbb{E}[r_i \mid s_i \in \Delta_k].
\end{equation}
AGILE advances only when the moving average on the new increment is high enough
and the cumulative performance does not fall below a floor:
\begin{equation}
    \overline{p}^{\mathrm{new}}_t \geq \tau^{\mathrm{new}}_k
    \quad \land \quad
    \overline{p}^{\mathrm{cum}}_t \geq \tau^{\mathrm{cum}}_k.
\end{equation}
This criterion explicitly asks whether the newly introduced region has become
learnable, which is important for diagnosing and improving generalization.

\paragraph{Simulation and real-world expansion.}
AGILE uses different adaptation mechanisms depending on the training setting.
In simulation, the policy is optimized with reinforcement learning within the
current stage until either the advancement criterion is satisfied or a maximum
number of stage steps is reached. The system then expands the environment reset
distribution to the next stage. In real-world training, AGILE can additionally
use a DAgger-style data correction loop: rollouts with low task success or poor
execution quality are selected for expert correction, and the corrected
trajectories are added back to the training set. Thus, simulation experiments
mainly validate automatic distribution expansion under RL, while real-world
experiments instantiate the same generalization-driven expansion principle
through active corrective data aggregation.

\subsection{Sampling Newly Added Configurations}
\label{subsec:mixing}

When a new stage is activated, the cumulative set can be much larger than the
newly added set. Uniform sampling over $\mathcal{A}_k$ may therefore under-sample
the new configurations. AGILE optionally uses a short-term mixture schedule:
\begin{equation}
    s \sim
    \lambda_t \, \mathrm{Unif}(\Delta_k)
    + (1-\lambda_t) \, \mathrm{Unif}(\mathcal{A}_{k-1}),
\end{equation}
where $\lambda_t$ can decay with the number of training steps spent in the
current stage. This gives the policy more updates on newly introduced
configurations while preserving rehearsal on previously learned regions.

\subsection{Policy Optimization}
\label{subsec:policy_optimization}

AGILE is independent of the specific policy optimizer. In our implementation,
we train the VLA policy with PPO-style actor-critic updates. Rollouts are
collected from IsaacLab environments using the active configuration set
$\mathcal{A}_k$. The environment returns standard episode rewards and success
signals, while the expansion controller consumes only aggregated success
statistics for stage advancement. Thus the learning objective remains unchanged;
AGILE changes the training distribution rather than the reward definition.

Let $\mathcal{D}_k$ denote rollouts collected at stage $k$. The policy is updated
by maximizing the clipped policy objective with value learning:
\begin{equation}
    \mathcal{L}(\theta) =
    \mathbb{E}_{\mathcal{D}_k}
    \left[
    \min
    \left(
    \rho_t(\theta) \hat{A}_t,
    \operatorname{clip}(\rho_t(\theta), 1-\epsilon, 1+\epsilon)\hat{A}_t
    \right)
    \right],
\end{equation}
where $\rho_t(\theta)$ is the policy likelihood ratio and $\hat{A}_t$ is the
advantage estimate. The same optimizer and model architecture can therefore be
used with or without AGILE.

\subsection{Implementation Details}
\label{subsec:implementation_details}

Each task configuration is stored as one JSONL record with a unique scenario id.
The stage manifest stores exact ids for each incremental and cumulative stage.
At runtime, the scheduler samples from the active ids and applies the selected
configuration during environment reset. Updating the active stage only changes
the scheduler's allowed id set; it does not recreate the IsaacLab environment.
This is critical for keeping scene switching efficient.

In the main 10-stage setting, we use 200 configurations with cumulative stage
sizes:
\begin{equation}
    (10, 20, 40, 60, 80, 100, 120, 140, 170, 200).
\end{equation}
The first stage therefore contains only 10 configurations, and later stages
gradually expand the training distribution. For larger-scale experiments, we
also use 19-stage variants over 10,025 configurations.

\begin{algorithm}[t]
\caption{AGILE: Adaptive Generalization via Iterative Learning and Expansion}
\label{alg:agile}
\begin{algorithmic}[1]
\REQUIRE Candidate configurations $\mathcal{S}$, embeddings $\{z_i\}$, stage sizes $\{n_k\}_{k=0}^{K}$, success thresholds $\{\tau_k\}$, hit counts $\{h_k\}$
\STATE Build stage manifest $\{\Delta_k, \mathcal{A}_k\}_{k=0}^{K}$ from world-model embeddings
\STATE Initialize policy $\pi_\theta$ and stage index $k \leftarrow 0$
\STATE Set active training configurations to $\mathcal{A}_0$
\WHILE{training is not finished}
    \STATE Sample configurations $s \sim \mathcal{A}_k$ and collect rollouts
    \STATE Update $\pi_\theta$ with policy optimization
    \STATE Estimate success rate on the current stage
    \IF{stage advancement criterion is satisfied and $k < K$}
        \STATE $k \leftarrow k + 1$
        \STATE Update active configurations to $\mathcal{A}_k$
        \STATE Reset stage-level success statistics
    \ENDIF
\ENDWHILE
\RETURN Trained policy $\pi_\theta$
\end{algorithmic}
\end{algorithm}
